\begin{figure}[!htbp]
\centering
\begin{adjustbox}{max width=\linewidth}
\begin{tikzpicture}[x=0.85cm]
  \foreach \s/\c [count=\i from 0] in {X/fsage, X/fsage, a/fpaper, a/fpaper, Y/fsand, Y/fsand, b/fpaper, b/fpaper, B/fgrey, B/fgrey} {
    \filldraw[fill=\c, draw=fgink, line width=0.5pt] (\i,0) rectangle (\i+1,0.85);
    \node[font=\small] at (\i+0.5,0.425) {$\s$};
  }
  \draw[thinline] (10,0) -- (10.6,0) (10,0.85) -- (10.6,0.85); \node at (10.9,0.425) {$\cdots$};
  \fill[fwine] (2.5,-0.15) -- (2.3,-0.55) -- (2.7,-0.55) -- cycle;
  \node[font=\footnotesize, text=fwine, below] at (2.5,-0.55) {head, state $q_0$};
  \node[lbl, anchor=west] at (11.4,0.425) {after two rounds on $a^4b^4$};
\end{tikzpicture}
\end{adjustbox}
\caption{A Turing machine for $a^n b^n$ repeatedly marks the leftmost $a$ as $X$, moves right to the leftmost
$b$ and marks it $Y$, then returns left. When no $a$ remains it checks that no $b$ remains, and accepts.}
\end{figure}
